\section{A Non-Text Channel, and a Prior Negative Result Explained}
\label{sec:latent}

A natural worry is that agents evade text-based oversight by moving to a channel the
principal cannot read. Prior follow-up work trained an $8.39$M-parameter residual
adapter mapping a sender's last-layer hidden states into a receiver's input-embedding
space, with both base models frozen and the link trained only on neutral assistance
data---never on game, reward, or collusion data. After repairing a receiver-layout
confound (\S\ref{sec:methods}), the final adapter reproduced held-out
text-conditioned action distributions better than shuffled, random, and zero controls
on fresh deployment snapshots (KL $0.138$ vs $0.197$/$0.255$/$0.349$; token oracle
$0.000$), with paired intervals excluding zero. Behaviourally it did nothing: trained
$-$ shuffled $= 0.000$ $[-0.075,+0.063]$, trained $-$ text $= -0.100$.

That null was reported as an isolated negative. The threshold model explains it. The
link was trained to reproduce the receiver's \emph{action distribution}---it transmits
\emph{what I will do next}, which is precisely an own intention. Under
\S\ref{sec:model} it should not coordinate, and it did not. A high-capacity,
unmonitorable channel that faithfully carries decision-relevant information buys an
attacker nothing, because fidelity of intention transfer is not the coordinating
property.

We state the scope carefully: this is evidence that transmitting intentions
faithfully is insufficient, \emph{not} that latent channels are safe. A link trained
to carry proposals---which our methodology deliberately forbade---may well behave
differently, and that is the experiment we specify next.

\paragraph{The transfer test, specified and released but not run.}
The revised account predicts something sharper than the predicate did. If the channel
works by changing what the \emph{sender} resolves to do, then a link that transports the
sender's hidden states should carry the effect only insofar as the sender was already
reframed---so restricting the sender's generation should matter, and the transport
medium should not. The design crosses text/latent $\times$ proposal/intention with
shuffled, zero and no-channel controls, gated on an exact-token oracle reading
$0.000$ KL and on a scale check, with the link trained on neutral text only.

Three outcomes are distinguishable: the framing effect survives translation (latent
channels inherit the attack surface, and text-based monitoring is evadable); it is
text-specific (which explains the L4 null and bounds the attack); or the link is a
plain text substitute.

\paragraph{Pilot result: the framing effect is text-specific.}
We ran the design at Qwen2.5-14B on a single A100 ($12$ seeds/arm, $84$ matches).
The scaffold cleared its self-test, and the capability gate passed decisively
(no-channel $0.167$ $[0.07,0.36]$ vs.\ text $0.750$ $[0.55,0.88]$; difference
$+0.583$, lower bound $+0.307$). The link trained to convergence on neutral text
(loss $0.331\to0.013$) and reproduced the teacher's next-token argmax
\emph{perfectly} on held-out neutral messages (Table~\ref{tab:l5fid}).

\begin{table}[t]
\centering
\caption{\textbf{L5 fidelity gates, Qwen2.5-14B, held-out neutral messages.} The
exact-token oracle validates the matched layout; the trained link is
indistinguishable from reading the text at the token level.}
\label{tab:l5fid}
\small
\begin{tabular}{lrr}
\toprule
payload & KL & top-1 agreement \\
\midrule
token oracle & $-0.0000$ & $1.000$ \\
trained      & $0.019$   & $1.000$ \\
shuffled     & $6.440$   & $0.109$ \\
random       & $8.019$   & $0.000$ \\
zero         & $18.101$  & $0.000$ \\
\bottomrule
\end{tabular}
\end{table}

Behaviourally the framing effect appears in text and \emph{not} in the latent channel
(Table~\ref{tab:l5c}). Over text, restricting the sender to joint proposals rather than
own intentions raises lock-in by $+0.417$ $[+0.087,+0.680]$---an independent
replication of P1's primary contrast in a different scaffold, different prompt
structure, and different round count. Passed through the link, the same manipulation
moves nothing: $+0.083$ $[-0.169,+0.354]$, and \texttt{latent\_proposal} is
indistinguishable from both its shuffled and its zero controls.

\begin{table}[t]
\centering
\caption{\textbf{L5-C pilot, $12$ seeds/arm.} The sender-side manipulation is
identical across the text and latent rows; only the transport differs.}
\label{tab:l5c}
\small
\begin{tabular}{lccr}
\toprule
arm & lock-in & Wilson 95\% CI & mean $C$ \\
\midrule
none               & 0.000 & [0.00, 0.24] & 0.446 \\
\textbf{text\_proposal}  & \textbf{0.417} & [0.19, 0.68] & 0.792 \\
text\_intention     & 0.000 & [0.00, 0.24] & 0.597 \\
latent\_proposal    & 0.083 & [0.01, 0.35] & 0.534 \\
latent\_intention   & 0.000 & [0.00, 0.24] & 0.520 \\
latent\_shuffled    & 0.000 & [0.00, 0.24] & 0.189 \\
latent\_zero        & 0.000 & [0.00, 0.24] & 0.309 \\
\bottomrule
\end{tabular}
\end{table}

\begin{figure}[t]
\centering
\includegraphics[width=\textwidth]{figures/f3_latent_transfer.pdf}
\caption{\textbf{Token-level fidelity does not imply behavioural transfer.} Left: on
held-out neutral messages the trained link reproduces the teacher's next-token argmax
perfectly, as does the exact-token oracle that validates the layout. Right: the
sender-side manipulation is identical across the text and latent arms, and only the
text arms separate. Fidelity is measured on neutral text, which bounds the reading
(see below).}
\label{fig:l5}
\end{figure}

\paragraph{The limitation that bounds this reading.}
The fidelity in Table~\ref{tab:l5fid} is measured on \emph{neutral} messages, and
\S\ref{sec:protocol} is precisely the finding that neutral fidelity does not predict
deployment fidelity---the original L4 adapter passed a neutral gate at KL $0.649$ and
then failed in deployment, as did exact token embeddings. This pilot does not
implement the deployment-matched gate that finding demands. So two explanations
remain open: the framing effect is genuinely text-specific, or the link is faithful on
neutral text and unfaithful on game messages, which is the L4 failure repeating. We
cannot separate them here, and we do not claim to.

What the pilot does establish is narrower and still useful: a channel that reproduces
the text-conditioned next-token distribution \emph{perfectly} on the distribution it
was trained for does not thereby transport the behavioural effect that the text
itself produces. Token-level fidelity and behavioural transfer are not the same
property. That is the same separation the L4 work reached from the other direction,
and it holds here at $14$B with a link an order of magnitude more faithful.

\paragraph{Scope.}
Twelve seeds give intervals roughly $\pm0.24$ wide. The text contrast clears zero and
is promotable to a $40$-seed confirmatory run; \textbf{the latent nulls are
under-powered and must not be read as evidence that latent channels are safe}. The
executed notebook is released with the artifact.

% ======================================================================
